\subsection{Measures defined by numerical densities}

Let \(g\) be a positive numerical function defined locally \(\mu\) -almost everywhere in \(T\) and locally \(\mu\) -integrable; the set of \(t\) such that \(g(t) = +\infty\) is then locally \(\mu\) -negligible, because \(g\varphi_{K}\) is \(\mu\) -integrable for every compact set \(K\) (Ch. IV, §2, No.~3, Prop. 7). Now let \(g'\) be a locally integrable function that is positive and finite, equal to \(g\) locally \(\mu\) -almost everywhere; set \(\lambda_t' = g'(t)\varepsilon_t\) . The mapping \(t \mapsto \lambda_t'\) of \(T\) into \(\mathcal{M}_+(T)\) is vaguely \(\mu\) -measurable and scalarly essentially integrable (or again, the pair \((I,g')\) , where I is the identity mapping of \(T\) , is \(\mu\) -adapted); the integral \(\nu = \int \lambda_t' d\mu(t)\) does not depend on the particular function \(g'\) , locally almost everywhere equal to \(g\) , used in the definition of the measures \(\lambda_t'\) . This measure \(\nu\) is determined by the condition

\[
\int f (t) d \nu (t) = \int f (t) g (t) d \mu (t) \quad \mathrm{for} f \in \mathcal {K} (\mathrm{T}).\tag{1}
\]

If now \(\theta\) is a complex measure, and if u is a complex function (or a function with values in \(\overline{R}\) ) defined locally \(\theta\) -almost everywhere and locally integrable for \(\theta\) , one can write

\[
\begin{array}{l} {u = g _ {1} - g _ {2} + i (g _ {3} - g _ {4})} \\ {\theta = \mu_ {1} - \mu_ {2} + i (\mu_ {3} - \mu_ {4})} \end{array}\tag{2}
\]

where \(\mu_{1} = (\mathcal{R}\theta)^{+}\) , \(\mu_{2} = (\mathcal{R}\theta)^{-}\) , \(\mu_{3} = (\mathcal{I}\theta)^{+}\) , \(\mu_{4} = (\mathcal{I}\theta)^{-}\) (Ch. III, §1, No.~5), and where \(g_{1}, g_{2}, g_{3}, g_{4}\) have the analogous meanings; since \(|u|\) is locally \(|\theta|\) -integrable, each of the positive functions \(g_{i}\) ( \(i = 1, 2, 3, 4\) ) is locally integrable for each positive measure \(\mu_{j}\) ( \(j = 1, 2, 3, 4\) ), so that the mapping

\[
f \mapsto \int f (t) u (t) d \theta (t)
\]

on \(\mathcal{K}(\mathrm{T})\) is a complex measure.

DEFINITION 2. --- Let \(\theta\) be a complex measure, and let \(u\) be a complex function (or a function with values in \(\overline{\mathbf{R}}\) ) defined locally \(\theta\) -almost everywhere and locally \(\theta\) -integrable. The complex measure \(f \mapsto \int f u d\theta\) on T is said to be the product of the measure \(\theta\) by the function \(u\) , or the measure with density \(u\) with respect to \(\theta\) , and is denoted \(u \cdot \theta\) .

Every complex measure that is the product of a positive measure \(\mu\) by a locally \(\mu\) -integrable function is called a measure with base \(\mu\) .

The relation \(\eta = u \cdot \theta\) is again, by convention, written

\[
d \eta (t) = u (t) d \theta (t).
\]

When u is everywhere defined and continuous, one recovers the definition given in Ch. III, §1, No.~4. It is clear that if \(u_{1}\) and \(u_{2}\) are locally \(\theta\) -integrable, then \((u_{1}+u_{2})\cdot\theta=u_{1}\cdot\theta+u_{2}\cdot\theta\) . Similarly, if \(\theta_{1}\) and \(\theta_{2}\) are two measures on T, and if u is a function locally integrable for \(\theta_{1}\) and \(\theta_{2}\) , then u is locally integrable for \(\theta_{1}+\theta_{2}\) and one has \(u\cdot(\theta_{1}+\theta_{2})=u\cdot\theta_{1}+u\cdot\theta_{2}\) .

We shall henceforth restrict ourselves to the case of functions defined everywhere; the extension to functions defined locally almost everywhere, which is always obvious, is left to the reader.

The following proposition permits, for the most part, reducing the case of complex measures to that of positive measures:

PROPOSITION 2. --- Let \(\theta\) be a complex measure, and \(u\) a locally \(\theta\) -integrable complex function; then

\[
\left| u \cdot \theta \right| = \left| u \right| \cdot \left| \theta \right|.\tag{3}
\]

We begin with an auxiliary result:

Lemma 1. --- Let \(\theta\) be a complex measure, and let \(f\) be an element of \(\overline{\mathcal{L}}_{\mathbf{C}}^{1}(\mathrm{T},\theta)\) . Then

\[
\langle | \theta |, | f | \rangle = \sup _ {c \in \mathcal {K} _ {1}} | \langle \theta , c f \rangle | = \sup _ {c \in \mathcal {B} _ {1}} | \langle \theta , c f \rangle |,\tag{4}
\]

where \(\mathcal{K}_{1}\) (resp. \(B_{1}\) ) denotes the set of complex functions c, continuous with compact support (resp. Borel), such that \(|c| \leqslant 1\) .

Let us first treat the case that \(f \in \mathcal{K}(\mathrm{T};\mathbf{C})\) . Obviously

\[
\sup _ {c \in \mathcal {K} _ {1}} | \langle \theta , c f \rangle | \leqslant \sup _ {c \in \mathcal {B} _ {1}} | \langle \theta , c f \rangle | \leqslant \langle | \theta |, | f | \rangle
\]

(Ch. IV, §4, No.~2, Prop. 2). On the other hand, let \(g\) be an element of \(\mathcal{K}(\mathrm{T};\mathbf{C})\) such that \(|g| \leqslant |f|\) ; \(g\) is the uniform limit of a sequence \((g_n)\) of elements of \(\mathcal{K}(\mathrm{T};\mathbf{C})\) whose supports are contained in the open set \(\mathrm{U}\) formed by the \(t\) such that \(f(t) \neq 0\) , and one may clearly suppose that \(|g_n| \leqslant |f|\) for every \(n\) . Set \(c_n(t) = g_n(t)/f(t)\) for \(t \in \mathrm{U}\) , \(c_n(t) = 0\) for \(t \notin \mathrm{U}\) ; then \(c_n \in \mathcal{K}_1\) , \(g = \lim_{n \to \infty} c_n f\) , therefore \(|\langle \theta, g \rangle| = \lim_{n \to \infty} |\langle \theta, c_n f \rangle|\) , and finally

\[
\sup _ {| g | \leqslant | f |, g \in \mathscr {K} (\mathrm{T}; \mathbf {C})} | \langle \theta , g \rangle | \leqslant \sup _ {c \in \mathscr {K} _ {1}} | \langle \theta , c f \rangle |.
\]

One concludes by observing that the first member of this inequality is equal to \(\langle|\theta|,|f|\rangle\) (Ch. III, §1, No.~6, formula (12)).

Next, denote by \(f\) an element of \(\overline{\mathcal{L}}_{\mathbf{C}}^{1}(\theta)\) , and let us show that (4) is again true: it suffices to verify that the three members of this relation depend continuously on \(f\) for the topology of \(\overline{\mathcal{L}}_{\mathbf{C}}^{1}(\theta)\) , since they coincide on the dense subspace \(\mathcal{K}(\mathrm{T};\mathbf{C})\) . This results at once from the following inequalities, where \(f\) and \(f'\) denote elements of \(\overline{\mathcal{L}}_{\mathbf{C}}^{1}(\theta)\) :

\[
| \langle | \theta |, | f | \rangle - \langle | \theta |, | f ^ {\prime} | \rangle | \leqslant \langle | \theta |, | f - f ^ {\prime} | \rangle = \overline {{\mathrm{N}}} _ {1} (f - f ^ {\prime})
\]

\[
| \langle \theta , c f \rangle - \langle \theta , c f ^ {\prime} \rangle | \leqslant \langle | \theta |, | c | | f - f ^ {\prime} | \rangle \leqslant \overline {{{\mathrm{N}}}} _ {1} (f - f ^ {\prime})
\]

for all \(c \in B_{1}\) . The lemma is thus established.

Passing to the proof of Proposition 2, let us apply the lemma to the function uf, where h belongs to \(\mathcal{K}_{+}(\mathrm{T})\) . This yields:

\[
\langle | \theta |, | u h | \rangle = \sup _ {c \in \mathscr {K} _ {1}} | \langle \theta , c u h \rangle | = \sup _ {c \in \mathscr {K} _ {1}} | \langle u \cdot \theta , c h \rangle | = \langle | u \cdot \theta |, h \rangle .
\]

However, the first member is also equal to

\[
\langle | \theta |, | u | h \rangle = \langle | u | \cdot | \theta |, h \rangle .
\]

The two measures \(|u| \cdot |\theta|\) and \(|u \cdot \theta|\) are therefore equal.

COROLLARY. --- Let \(g_{1}\) and \(g_{2}\) be two locally \(\mu\) -integrable numerical functions; then

\[
\inf (g _ {1} \cdot \mu , g _ {2} \cdot \mu) = \inf (g _ {1}, g _ {2}) \cdot \mu ; \quad \sup (g _ {1} \cdot \mu , g _ {2} \cdot \mu) = \sup (g _ {1}, g _ {2}) \cdot \mu .
\]

In particular, if \(g\) is a locally \(\mu\) -integrable numerical function, then

\[
(g \cdot \mu) ^ {+} = g ^ {+} \cdot \mu ; (g \cdot \mu) ^ {-} = g ^ {-} \cdot \mu .
\]

This follows at once from Prop. 2 and the formulas (6) of Ch. II, §1, No.~1.

